\documentclass[12pt,a4paper]{article}

\usepackage[german]{babel}
\usepackage[utf8]{inputenc}
\usepackage[T1]{fontenc}
\usepackage{enumerate}
\usepackage[absolute]{textpos}
\usepackage[a4paper, top=2cm, left=2cm, bottom=2cm, right=2cm]{geometry}
\usepackage{parskip}

\usepackage{graphicx}

\begin{document}
\thispagestyle{empty}

\begin{textblock*}{\paperwidth}(1.6cm,0.48cm) %{width}(horizontal, vertical)
\includegraphics[scale=0.67]{Logo_Mathematik_CMYK.pdf}
\end{textblock*}

\vspace*{1.5cm}

\begin{center}
\Huge{\textbf{Informationen über den Prüfungsablauf}}\\
\large{Mathematik-Olympiade 2019}
\end{center}

\bigskip

Die untenstehenden Informationen regeln den Ablauf der Prüfungen im Rahmen der Mathematik-Olympiade.

\section{Erlaubte Hilfsmittel}
\begin{itemize}
\item Erlaubte Hilfsmittel sind Stifte und Geometriewerkzeuge (Zirkel, Lineal, Winkelmass). 
\item Der Gebrauch von Taschenrechnern und Formelsammlungen ist verboten. Eigenes Papier zu verwenden ist nicht gestattet.
\item Es ist erlaubt, Essen und Trinken an die Prüfung mitzunehmen. Aus Rücksicht auf die anderen Teilnehmenden soll auf laute oder geruchsintensive Nahrung verzichtet werden.
\end{itemize}

\section{Ablauf der Prüfung}
\begin{itemize}

\item Die Lösungen und Skizzen müssen auf Papier, das von der Prüfungsaufsicht ausgehändigt wird, aufgeschrieben werden. Es ist möglich, während der Prüfung zusätzliche Blätter zu verlangen.

\item Es können jederzeit Fragen gestellt werden, welche die Aufgabenstellung betreffen. Alle Fragen zu einer Aufgabe müssen schriftlich auf einem separaten Blatt Papier einer Aufsichtsperson abgegeben werden.

\item Es ist nur mit Erlaubnis einer Aufsichtsperson gestattet, während der Prüfung auf die Toilette zu gehen.

\end{itemize}

\section{Formales und Korrektur}
\begin{itemize}
\item Es darf nichts auf die Rückseite der Blätter geschrieben werden. Für verschiedene Aufgaben müssen verschiedene Blätter benutzt werden.
\item Auf jedem Blatt müssen Name, Aufgabennummer, Seitennummer und die totale Anzahl Seiten für diese Aufgabe ersichtlich sein, z.B: "Hermann Schwarz, G1, S.2/4"
\item Alle Schritte der Lösungen müssen begründet sein. Für alle verwendeten mathematischen Resultate muss eine passende Referenz angegeben werden.
\item Nach Erhalt der korrigierten Prüfung ist es möglich, innert der festgelegten Frist Erläuterungen für die Korrektur zu verlangen. Falls angemessen, kann die Korrektur noch einmal überprüft werden.
\end{itemize}
\end{document} 